\documentclass[10pt,a4paper]{amsart}
\usepackage{amsmath,amssymb,amsthm,graphicx}
\usepackage{verbatim}
\newenvironment{reference}{\comment}{\endcomment}
%\newenvironment{reference}{}{}
\newenvironment{slogan}{\comment}{\endcomment}
\newenvironment{history}{\comment}{\endcomment}

% For commutative diagrams we use Xy-pic
\usepackage[all]{xy}

% We use 2cell for 2-commutative diagrams.
\xyoption{2cell}
% Theorem environments.
%
\theoremstyle{plain}
\newtheorem{theorem}[subsection]{Theorem}
\newtheorem{proposition}[subsection]{Proposition}
\newtheorem{lemma}[subsection]{Lemma}

\theoremstyle{definition}
\newtheorem{definition}[subsection]{Definition}
\newtheorem{example}[subsection]{Example}
\newtheorem{exercise}[subsection]{Exercise}
\newtheorem{situation}[subsection]{Situation}

\theoremstyle{remark}
\newtheorem{remark}[subsection]{Remark}
\newtheorem{remarks}[subsection]{Remarks}

\numberwithin{equation}{subsection}

% Macros
%
\def\lim{\mathop{\mathrm{lim}}\nolimits}
\def\colim{\mathop{\mathrm{colim}}\nolimits}
\def\Spec{\mathop{\mathrm{Spec}}}
\def\Hom{\mathop{\mathrm{Hom}}\nolimits}
\def\Ext{\mathop{\mathrm{Ext}}\nolimits}
\def\SheafHom{\mathop{\mathcal{H}\!\mathit{om}}\nolimits}
\def\SheafExt{\mathop{\mathcal{E}\!\mathit{xt}}\nolimits}
\def\Sch{\mathit{Sch}}
\def\Mor{\mathop{\mathrm{Mor}}\nolimits}
\def\Ob{\mathop{\mathrm{Ob}}\nolimits}
\def\Sh{\mathop{\mathit{Sh}}\nolimits}
\def\NL{\mathop{N\!L}\nolimits}
\def\CH{\mathop{\mathrm{CH}}\nolimits}
\def\proetale{{pro\text{-}\acute{e}tale}}
\def\etale{{\acute{e}tale}}
\def\QCoh{\mathit{QCoh}}
\def\Ker{\mathop{\mathrm{Ker}}}
\def\Im{\mathop{\mathrm{Im}}}
\def\Coker{\mathop{\mathrm{Coker}}}
\def\Coim{\mathop{\mathrm{Coim}}}

% Boxtimes
%
\DeclareMathSymbol{\boxtimes}{\mathbin}{AMSa}{"02}

%
% Macros for moduli stacks/spaces
%
\def\QCohstack{\mathcal{QC}\!\mathit{oh}}
\def\Cohstack{\mathcal{C}\!\mathit{oh}}
\def\Spacesstack{\mathcal{S}\!\mathit{paces}}
\def\Quotfunctor{\mathrm{Quot}}
\def\Hilbfunctor{\mathrm{Hilb}}
\def\Curvesstack{\mathcal{C}\!\mathit{urves}}
\def\Polarizedstack{\mathcal{P}\!\mathit{olarized}}
\def\Complexesstack{\mathcal{C}\!\mathit{omplexes}}
% \Pic is the operator that assigns to X its picard group, usage \Pic(X)
% \Picardstack_{X/B} denotes the Picard stack of X over B
% \Picardfunctor_{X/B} denotes the Picard functor of X over B
\def\Pic{\mathop{\mathrm{Pic}}\nolimits}
\def\Picardstack{\mathcal{P}\!\mathit{ic}}
\def\Picardfunctor{\mathrm{Pic}}
\def\Deformationcategory{\mathcal{D}\!\mathit{ef}}
\usepackage{cleveref}
\crefname{originalreference}{original reference}{original references}
\Crefname{originalreference}{Original reference}{Original references}
\makeatletter
\newlabel{lemma-etale-gives-rig-etale}{{0AR2}{source}{}{sourcepointer-1}{}}
\newlabel{lemma-etale-gives-rig-etale@cref}{{[originalreference][0][]0AR2}{[1][0][]source}{}{}{}}
\makeatother
\crefname{originalsection}{original section}{original sections}
\Crefname{originalsection}{Original section}{Original sections}
\makeatletter
\newlabel{section-completion-functor}{{0GDB}{source}{}{sourcepointer-2}{}}
\newlabel{section-completion-functor@cref}{{[originalsection][0][]0GDB}{[1][0][]source}{}{}{}}
\makeatother
\title{Functoriality of the completion functor under base change}
\author{The Stacks Project Authors}
\date{stacks-project:restricted.tex:local-snapshot-20261012; GNU FDL 1.2 or later}
\begin{document}
\maketitle
\noindent\textit{Editorial note.} All delivered mathematical text is the author's own, copied from the pinned Stacks Project chapter file restricted.tex; no statement or proof step was rewritten, reordered or removed.\par\medskip
\noindent\textit{Editorial note.} The counted proof is the author's complete argument: the counted statement is the author lemma environment labelled lemma-functoriality-completion-functor; the proof environment immediately following it is its own proof; the context span delivers the section situation and the completion functor equation that 'in the situation above' refers to\par\medskip
\noindent\textit{Editorial note.} The label lemma-functoriality-completion-functor names the counted statement, and the proof environment immediately following it in the source is delivered verbatim as its proof.\par\medskip
\noindent\textit{Editorial note.} The counted statement is carried by the author's own lemma environment in the chapter source.\par\medskip
\noindent\textit{Allowed standard prerequisites (editorial).}\begin{itemize}
\item the section-algebraization setup: locally Noetherian X, closed T, U, and the completion functor (h : Y to X) mapping to Y\_\{/T\}, together with the category definitions (delivered as context)
\item the completion functor commutes with fibre products (formal-spaces-section-completion-functor)
\end{itemize}\medskip
\ifdefined\hypertarget\hypertarget{sourcepointer-1}{}\fi
\noindent\textit{Source reference.} Original reference 0AR2: sources/stacks-restricted/pinned-original-source.tex, source lines 6186--6186. Prerequisite pointer invoked at 6878; it names a standard fact of the same chapter that the counted argument uses, and every step of the counted proof built on it is retained verbatim in the delivered body.\par\medskip
\ifdefined\hypertarget\hypertarget{sourcepointer-2}{}\fi
\noindent\textit{Source reference.} Original section 0GDB: sources/stacks-restricted/pinned-original-source.tex, source lines 5953--5953. Prerequisite pointer invoked at 6832; it names a standard fact of the same chapter that the counted argument uses, and every step of the counted proof built on it is retained verbatim in the delivered body.\par\medskip
\medskip\noindent
The notation in this section will agree with the notation in
Section \ref{section-completion-functor} except our algebraic
spaces and formal algebraic spaces will be locally Noetherian.

\medskip\noindent
Thus, we first fix a base scheme $S$. All rings, topological rings,
schemes, algebraic spaces, and formal algebraic spaces and
morphisms between these will be over $S$. Next, we fix a
locally Noetherian algebraic space $X$ and a closed subset $T \subset |X|$.
We denote $U \subset X$ be the open subspace with $|U| = |X| \setminus T$.
Picture
$$
U \to X \quad |X| = |U| \amalg T
$$
Given a morphism of algebraic spaces $f : X' \to X$, we will use the notation
$U' = f^{-1}U$, $T' = |f|^{-1}(T)$, and $f_{/T} : X'_{/T'} \to X_{/T}$ as in
Section \ref{section-completion-functor}. We will sometimes write
$X'_{/T}$ in stead of $X'_{/T'}$ and more generally for a morphism
$a : X' \to X''$ of algebraic spaces over $X$ we will denote
$a_{/T} : X'_{/T} \to X''_{/T}$ the induced morphism of formal
algebraic spaces obtained by completing the morphism $a$ along
the inverse images of $T$ in $X'$ and $X''$.

\medskip\noindent
Given this setup we will consider the functor
\begin{equation}
\label{equation-completion-functor}
\left\{
\begin{matrix}
\text{morphisms of algebraic spaces}\\
f : X' \to X\text{ which are locally}\\
\text{of finite type and such that}\\
U' \to U\text{ is an isomorphism}
\end{matrix}
\right\}
\longrightarrow
\left\{
\begin{matrix}
\text{morphisms }g : W \to X_{/T}\\
\text{of formal algebraic spaces}\\
\text{with }W\text{ locally Noetherian}\\
\text{and }g\text{ rig-\'etale}
\end{matrix}
\right\}
\end{equation}
sending $f : X' \to X$ to $f_{/T} : X'_{/T'} \to X_{/T}$.
This makes sense because $f_{/T}$ is rig-\'etale by
Lemma \ref{lemma-etale-gives-rig-etale}.

\begin{lemma}
\label{lemma-functoriality-completion-functor}
In the situation above, let $X_1 \to X$ be a morphism of algebraic
spaces with $X_1$ locally Noetherian. Denote $T_1 \subset |X_1|$
the inverse image of $T$ and $U_1 \subset X_1$ the inverse image of $U$.
We denote
\begin{enumerate}
\item $\mathcal{C}_{X, T}$ the category whose objects are
morphisms of algebraic spaces $f : X' \to X$ which are locally
of finite type and such that $U' = f^{-1}U \to U$ is an isomorphism,
\item $\mathcal{C}_{X_1, T_1}$ the category whose objects are
morphisms of algebraic spaces $f_1 : X_1' \to X_1$ which are locally
of finite type and such that $f_1^{-1}U_1 \to U_1$ is an isomorphism,
\item $\mathcal{C}_{X_{/T}}$ the category whose objects are
morphisms $g : W \to X_{/T}$ of formal algebraic spaces
with $W$ locally Noetherian and $g$ rig-\'etale,
\item $\mathcal{C}_{X_{1, /T_1}}$ the category whose objects are
morphisms $g_1 : W_1 \to X_{1, /T_1}$ of formal algebraic spaces
with $W_1$ locally Noetherian and $g_1$ rig-\'etale.
\end{enumerate}
Then the diagram
$$
\xymatrix{
\mathcal{C}_{X, T} \ar[d] \ar[r] &
\mathcal{C}_{X_{/T}} \ar[d] \\
\mathcal{C}_{X_1, T_1} \ar[r] &
\mathcal{C}_{X_{1, /T_1}}
}
$$
is commutative where the horizontal arrows are given by
(\ref{equation-completion-functor})
and the vertical arrows by base change along
$X_1 \to X$ and along $X_{1, /T_1} \to X_{/T}$.
\end{lemma}

\begin{proof}
This follows immediately from the fact that the completion functor
$(h : Y \to X) \mapsto Y_{/T} = Y_{/|h|^{-1}T}$
on the category of algebraic spaces over $X$
commutes with fibre products.
\end{proof}

\end{document}
